\documentclass{beamer}
\usepackage{tikz,amsmath,hyperref,graphicx,stackrel,animate,media9}
\usetikzlibrary{positioning,shadows,arrows,shapes,calc}
\newcommand{\argmax}{\operatornamewithlimits{argmax}}
\newcommand{\argmin}{\operatornamewithlimits{argmin}}
\mode<presentation>{\usetheme{Frankfurt}}
\AtBeginSection[]
{
  \begin{frame}<beamer>
    \frametitle{Outline}
    \tableofcontents[currentsection,currentsubsection]
  \end{frame}
}
\title{Lecture 28: Z Transform}
\author{Mark Hasegawa-Johnson}
\date{ECE 401: Signal and Image Analysis}  
\begin{document}

% Title
\begin{frame}
  \maketitle
\end{frame}

% Title
\begin{frame}
  \tableofcontents
\end{frame}

%%%%%%%%%%%%%%%%%%%%%%%%%%%%%%%%%%%%%%%%%%%%
\section[Review]{Review: Frequency Response}
\setcounter{subsection}{1}

\begin{frame}
  \frametitle{Review: Frequency Response}
  \begin{itemize}
  \item {\bf Tones in $\rightarrow$ Tones out}
    \begin{align*}
      x[n]=e^{j\omega n} &\rightarrow y[n]=H(\omega)e^{j\omega n}\\
      x[n]=\cos\left(\omega n\right)
      &\rightarrow y[n]=|H(\omega)|\cos\left(\omega n+\angle H(\omega)\right)\\
      x[n]=A\cos\left(\omega n+\theta\right)
      &\rightarrow y[n]=A|H(\omega)|\cos\left(\omega n+\theta+\angle H(\omega)\right)
    \end{align*}
  \item where the {\bf Frequency Response} is given by
    \[
    H(\omega) = \sum_m h[m]e^{-j\omega m}
    \]
  \end{itemize}
\end{frame}  

\begin{frame}
  \frametitle{Linearity and Time-Shift Properties}

  The linearity property of the DTFT says that
  \[
  z[n]= ax[n]+by[n]\leftrightarrow Z(\omega) = aX(\omega)+bY(\omega).
  \]
  The time-shift property says that
  \[
  z[n]= x[n-n_0] \leftrightarrow Z(\omega) = e^{-j\omega n_0} X(\omega)
  \]
\end{frame}

%%%%%%%%%%%%%%%%%%%%%%%%%%%%%%%%%%%%%%%%%%%%
\section{Difference Equations}
\setcounter{subsection}{1}

\begin{frame}
  \frametitle{Difference Equation}

  A {\bf difference equation} is an equation in terms of time-shifted
  values of multiple signals.  For example,
  \[
  y[n] = x[n] - 2x[n-1] + 2x[n-2]
  \]
  By combining the linearity and time-shift properties of the DTFT, we
  can translate the whole difference equation into the frequency
  domain as
  \[
  Y(\omega)  = X(\omega) -2e^{-j\omega} X(\omega) + 2 e^{-2j\omega}X(\omega)
  \]
\end{frame}

\begin{frame}
  \frametitle{Impulse Response}

  A difference equation implements a discrete-time filter.  Therefore,
  it has an {\bf impulse response}.  You can find the impulse response
  by just putting in an impulse, $x[n]=\delta[n]$, and seeing how it
  responds.  Whatever value of $y[n]$ that comes out of the filter is
  the impulse response:
  \begin{align*}
    h[n] &= \delta[n]-2\delta[n-1]+2\delta[n-2]=\begin{cases}
    1 & n=0\\-2 & n=1\\2 & n=2\\0 & \mbox{otherwise}\end{cases}
  \end{align*}
  If you wanted to use {\tt np.convolve} to implement this filter, you
  now know what impulse response to use.
  \begin{align*}
    y[n] &= x[n] -2x[n-1]+2x[n-2]\\
    & = \sum_{m=0}^2 h[m]x[n-m]
  \end{align*}
\end{frame}  

\begin{frame}
  \frametitle{Frequency Response}

  The {\bf frequency response} of a filter is the function $H(\omega)$ such that
  $Y(\omega)=H(\omega)X(\omega)$, or in other words,
  \[
  H(\omega) = \frac{Y(\omega)}{X(\omega)}
  \]
  We can get this from the DTFT of the difference equation:
  \[
  Y(\omega)  = \left(1-2e^{-j\omega}+2e^{-2j\omega}\right) X(\omega)
  \]
  \[
  H(\omega) = \frac{Y(\omega)}{X(\omega)} = \left(1-2e^{-j\omega}+2e^{-2j\omega}\right)
  \]
\end{frame}

\begin{frame}
  \frametitle{Frequency Response is DTFT of the Impulse Response}

  The frequency response is the DTFT of the impulse response.
  \begin{align*}
    h[n] &= \delta[n]-2\delta[n-1]+2\delta[n-2]\\
    H(\omega) &= 1-2e^{-j\omega}+2e^{-2j\omega}
  \end{align*}
  It sounds like an accident, that frequency response is DTFT of the
  impulse response.  But actually, it's because the method for
  computing frequency response and the method for computing impulse
  response do the same two steps, in the opposite order (see next slide).
\end{frame}

\begin{frame}
  \begin{block}{How to compute the frequency response}
    \begin{enumerate}
    \item Take the DTFT of every term:
      \[
      Y(\omega)= X(\omega) -2e^{-j\omega}X(\omega)+2e^{-2j\omega}X(\omega)
      \]
    \item Divide by $X(\omega)$.
      \[
      H(\omega)= 1-2e^{-j\omega} +2e^{-2j\omega}
      \]
    \end{enumerate}
  \end{block}

  \begin{block}{How to compute the DTFT of the impulse response}
    \begin{enumerate}
    \item Replace $x[n]$ by $\delta[n]$:
      \[
      h[n]= \delta[n] -2\delta[n-1] +2\delta[n-2]
      \]      
    \item Take the DTFT:
      \[
      H(\omega)= 1-2e^{-j\omega} + 2e^{-2j\omega}
      \]      
    \end{enumerate}
  \end{block}
\end{frame}


%%%%%%%%%%%%%%%%%%%%%%%%%%%%%%%%%%%%%%%%%%%%
\section{The Z Transform and the Transfer Function}
\setcounter{subsection}{1}

\begin{frame}
  \frametitle{Z: a Frequency Variable for Time Shifts}

  If we're going to be working a lot with delays, instead of pure
  tones, then it helps to change our frequency variable.
  Until now, we've been working in terms of $\omega$, the frequency of the pure tone:
  \[
  X(\omega) = \sum_{n=-\infty}^\infty x[n]e^{-j\omega n}
  \]
  A unit delay, $\delta[n-1]$, has the DTFT $e^{-j\omega}$.
  Just to reduce the notation a little, let's define the basic unit to be
  $z=e^{j\omega}$, which is the transform of $\delta[n+1]$ (a unit advance).  Then we  get:
  \[
  X(z)   = \sum_{n=-\infty}^\infty x[n]z^{-n}
  \]
\end{frame}

\begin{frame}
  \frametitle{Main Use of Z Transform: Difference Equations}
  
  The main purpose of the Z transform, for now, is just so that we have less to write.
  Instead of transforming
  \[
  y[n] = x[n] - 2x[n-1] + 2x[n-2]
  \]
  to get
  \[
  Y(\omega) = \left(1-2e^{-j\omega}+2e^{-2j\omega}\right)X(\omega)
  \]
  Now we can just write
  \[
  Y(z) = \left(1-2z^{-1}+2z^{-2}\right) X(z)
  \]
\end{frame}

\begin{frame}
  \frametitle{Transfer function $H(z)$ = $z$-transform version of frequency response}
  
  \begin{itemize}
  \item You can find the \textbf{frequency response} of a system by
    dividing $Y(\omega)/X(\omega)$:
    \[
    H(\omega)=\frac{Y(\omega)}{X(\omega)} = 1-2e^{-j\omega}+2e^{-2j\omega}
    \]
  \item By dividing $Y(z)/X(z)$, we get a function called the
    \textbf{transfer function}.  
    \[
    H(z) = \frac{Y(z)}{X(z)} = 1-2z^{-1}+2z^{-2}
    \]
  \item \textbf{frequency response}=\textbf{transfer function evaluated at} $z=e^{j\omega}$.
  \end{itemize}
\end{frame}

\begin{frame}
  \frametitle{Main Use of Z Transform: Difference Equations}
  
  The longer the difference equation, the more you will appreciate
  writing $z$ instead of $e^{j\omega}$.
  \begin{align*}
    y[n] &= 0.2x[n+3]+0.3x[n+2]+0.5x[n+1]\\
    &-0.5x[n-1]-0.3x[n-2]-0.2x[n-2]\\
  \end{align*}
  \[
  Y(z) = \left(0.2z^{3}+0.3z^{2}+0.5z^{1}-0.5z^{-1}-0.3z^{-2}-0.2z^{-3}\right)X(z)
  \]
\end{frame}


%%%%%%%%%%%%%%%%%%%%%%%%%%%%%%%%%%%%%%%%%%%%
\section[Autoregressive]{Autoregressive Difference Equations}
\setcounter{subsection}{1}

\begin{frame}
  \frametitle{Autoregressive Filter}

  Today, instead of filters with this form:
  \begin{displaymath}
    y[n] = x[n]-x[n-1]
  \end{displaymath}
  \ldots we want to start studying filters of this form:
  \begin{displaymath}
    y[n] = x[n]-y[n-1]
  \end{displaymath}
  \begin{itemize}
  \item This is called {\bf autoregressive}, meaning that $y[n]$ depends on
    past values (regressive) of itself (auto).
  \item This is also an {\bf infinite impulse response} (IIR) filter, because
    the impulse response ($h[n]$) is infinitely long.
  \end{itemize}
\end{frame}
    
\begin{frame}
  \frametitle{Autoregressive Difference Equations}

  An {\bf autoregressive} filter is one in which the output, $y[n]$,
  depends on past values of itself ({\bf auto}=self, {\bf
    regress}=go back).  For example,
  \[
  y[n] = x[n] + 0.3x[n-1] + 0.8 y[n-1]
  \]
  
\end{frame}

\begin{frame}
  \frametitle{FIR vs. IIR Filters}

  \begin{itemize}
  \item A difference equation with only feedforward terms is always a
    {\bf finite impulse response (FIR)} filter, because $h[n]$ has
    finite length.
  \item An autoregressive filter (with feedback terms) is usually
    \textbf{infinite impulse response (IIR)}, because $h[n]$
    is infinitely long.
  \end{itemize}
\end{frame}

\begin{frame}
  \frametitle{Causal and Anti-Causal Filters}

  \begin{itemize}
  \item If the outputs of a filter depend only on {\bf current and
    past} values of the input, then the filter is said to be {\bf
    causal}.  An example is
    \[
    y[n] = x[n] + 0.3x[n-1] + 0.8 y[n-1]
    \]
  \item Any system that depends on both past and future is called
    \textbf{non-causal}.  An example is:
    \[
    y[n] = \frac{1}{3}\left(x[n-1]+x[n]+x[n+1]\right)
    \]
  \item \textbf{Non-causal FIR} filters are easy to analyze, but
    \textbf{non-causal IIR} filters sometimes can't even be implemented.
    Therefore \textbf{all IIR filter in this class will be causal}.
  \end{itemize}
  
\end{frame}

\begin{frame}
  \frametitle{Autoregressive Difference Equations}

  We can find the frequency response by taking the Z transform or DTFT
  of each term in the equation:
  \begin{align*}
    y[n] &= x[n] + 0.3 x[n-1] +  0.8 y[n-1]\\
    Y(z) &= X(z) + 0.3z^{-1}X(z) + 0.8  z^{-1}Y(z)\\
    \left(1-0.8z^{-1}\right)Y(z) &= \left(1+0.3z^{-1}\right)X(z)\\
    H(z) &=\frac{Y(z)}{X(z)} = \frac{1+0.3z^{-1}}{1-0.8 z^{-1}}\\
    H(\omega) &= \left.H(z)\right|_{z=e^{j\omega}} = \frac{1+0.3e^{-j\omega}}{1-0.8 e^{-j\omega}}
  \end{align*}
\end{frame}

\begin{frame}
  \frametitle{Frequency Response}

  Here is the frequency response of this filter, plotted using {\tt
    np.abs((1+0.3*np.exp(-1j*omega))/(1-0.8*np.exp(-1j*omega)))}
  
  \centerline{\includegraphics[height=2.5in]{exp/intro.png}}
\end{frame}

\begin{frame}
  \frametitle{General form of an FIR filter}
  Every FIR filter is just a convolution with a finite-length $h[n]$:
  \begin{align*}
    y[n] &= \sum_{k=0}^{M} b_k x[n-k] = \sum_k h[k]x[n-k]\\
    h[k] &=\begin{cases}
    b_k & 0\le k\le M\\
    0 &\mbox{otherwise}
    \end{cases}
  \end{align*}
  This filter has an impulse response ($h[n]$) that is $M+1$ samples
  long.
  \begin{itemize}
  \item The $b_k$'s are called {\bf feedforward}
    coefficients, because they feed $x[n]$ forward into $y[n]$.
  \end{itemize}
\end{frame}

\begin{frame}
  \frametitle{General form of an IIR filter}
  For an IIR filter, $h[n]$ has infinite length, so writing out the convolution is
  inconvenient.  Instead, we usually write the \textbf{difference equation:}
  \[
  \sum_{\ell=0}^N a_\ell y[n-\ell] = \sum_{k=0}^{M} b_k x[n-k]
  \]
  \begin{itemize}
  \item The $a_\ell$'s are caled {\bf feedback} coefficients, because
    they feed $y[n]$ back into itself.
  \item Can we find $h[n]$ in terms of $a_\ell$ and $b_k$?
  \end{itemize}
\end{frame}

%%%%%%%%%%%%%%%%%%%%%%%%%%%%%%%%%%%%%%%%%%%%
\section[First-Order]{Impulse Response and Transfer Function of a First-Order Autoregressive Filter}
\setcounter{subsection}{1}

\begin{frame}
  \frametitle{First-Order Feedback-Only Filter}

  Let's find the general form of $h[n]$, for the simplest possible
  autoregressive filter: a filter with one feedback term, and no
  feedforward terms, like this:
  \[
  y[n] = x[n] + ay[n-1],
  \]
  where $a$ is any constant (positive, negative, real, or complex).
\end{frame}

\begin{frame}
  \frametitle{Impulse Response of a First-Order Filter}

  We can find the impulse response by putting in $x[n]=\delta[n]$, and
  getting out $y[n]=h[n]$:
  \[
  h[n] = \delta[n] + ah[n-1].
  \]
  Recursive computation gives
  \begin{align*}
    h[0] &= 1 \\
    h[1] &= a\\
    h[2] &= a^2\\
     & \vdots\\
    h[n] &= a^nu[n]
  \end{align*}
  where we use the notation $u[n]$ to mean the ``unit step function,''
  \[u[n] = \begin{cases}1& n\ge 0\\0 & n<0\end{cases}\]
\end{frame}

\begin{frame}
  \frametitle{Impulse Response of Stable First-Order Filters}

  The coefficient, $a$, can be positive, negative, or even complex.
  If $a$ is complex, then $h[n]$ is also complex-valued.
  \centerline{\includegraphics[height=2.5in]{exp/iir_stable.png}}

\end{frame}

\begin{frame}
  \frametitle{Impulse Response of Unstable First-Order Filters}

  If $|a|>1$, then the impulse response grows exponentially.  If
  $|a|=1$, then the impulse response never dies away.  In either case,
  we say the filter is ``unstable.''
  \centerline{\includegraphics[height=2.5in]{exp/iir_unstable.png}}

\end{frame}

\begin{frame}
  \frametitle{Instability}

  \begin{itemize}
  \item A {\bf stable} filter is one that always generates finite
    outputs ($|y[n]|$ finite) for every possible finite input
    ($|x[n]|$ finite).
  \item An {\bf unstable} filter is one that, at least sometimes,
    generates infinite outputs, even if the input is finite.
  \item A first-order IIR filter is stable if and only if $|a|<1$.
  \end{itemize}
\end{frame}


\begin{frame}
  \frametitle{Frequency Response of a First-Order Filter}

  If the filter is stable ($|a|<1$), then 
  we can find the frequency response by taking the Z transform of $h[n]$:
  \begin{displaymath}
    h[n] = \begin{cases}a^n & n\ge 0\\
      0 &n < 0
    \end{cases}
  \end{displaymath}
  \begin{align*}
    H(z) &= \sum_{n=-\infty}^\infty h[n]z^{-n}\\
    &= \sum_{n=0}^\infty \left(az^{-1}\right)^n\\
    &= \begin{cases}
      \frac{1}{1-az^{-1}} & |a| < 1\\
      \infty & |a| \ge 1
    \end{cases}\\
    H(\omega) &= 
    &= \begin{cases}
      \frac{1}{1-ae^{-j\omega}} & |a| < 1\\
      \infty & |a| \ge 1
    \end{cases}\\
  \end{align*}
\end{frame}

\begin{frame}
  \frametitle{Frequency Response of a First-Order Filter}

  If the filter is stable ($|a|<1$), then 
  we can also find the frequency response by taking the Z transform of each
  term in the filter equation:
  \begin{align*}
  y[n] &= x[n] - ay[n-1],\\
  Y(z) &= X(z)-az^{-1} Y(z),\\
  H(z) &= \frac{Y(z)}{X(z)}=\frac{1}{1-az^{-1}}\\
  H(\omega) &= \frac{Y(\omega)}{X(\omega)}=\frac{1}{1-ae^{-j\omega}}
  \end{align*}
  That {\bf looks} like it works even if $|a|\ge 1$, but it's a lie.
  If $|a|\ge 1$, then when you put a pure tone as input, you might get
  $y[n]=\infty$ as output instead of $y[n]=H(\omega)x[n]$.
\end{frame}

\begin{frame}
  \frametitle{Frequency Response of a  First-Order Filter}
  \[
  H(\omega) = \frac{1}{1-ae^{-j\omega}}~~~\mbox{iff}~|a|<1
  \]
  
  \centerline{\includegraphics[width=4.5in]{exp/iir_freqresponse.png}}
\end{frame}

\begin{frame}
  \frametitle{Try the quiz!}

  Try the quiz!
\end{frame}

%%%%%%%%%%%%%%%%%%%%%%%%%%%%%%%%%%%%%%%%%%%%
\section[Summary]{Summary}
\setcounter{subsection}{1}

\begin{frame}
  \frametitle{Summary: Z Transform}
  \begin{itemize}
  \item A {\bf difference equation} is an equation in terms of
    time-shifted copies of $x[n]$ and/or $y[n]$.
  \item We have less to write if we use a new frequency variable,
    $z=e^{j\omega}$.  This leads us to the \textbf{Z transform}:
    \[
    X(z) = \sum_{n=-\infty}^\infty x[n]z^{-n}
    \]
  \item We can find the \textbf{transfer function}
    $H(z)=Y(z)/X(z)$ by taking the Z transform of each term of
    the difference equation.  This will result in the ratio of two polynomials in $z$.
  \item We can find the \textbf{frequency response} by
    setting $z=e^{j\omega}$.
  \end{itemize}
\end{frame}

\begin{frame}
  \frametitle{Summary: Autoregressive Filter}
  \begin{itemize}
  \item An {\bf autoregressive filter} is a filter whose current output,
    $y[n]$, depends on  past values of the output.
  \item An autoregressive filter is usually {\bf infinite impulse response (IIR)},
    because $h[n]$ has infinite length.
  \item A filter with only feedforward coefficients, and no feedback coefficients, is called
    {\bf finite impulse response (FIR)}, because $h[n]$ has finite length (its length is
    just the number of feedforward terms in the difference equation).
  \item The first-order, feedback-only autoregressive filter has this
    impulse response and frequency response:
    \[
    h[n]=a^n u[n] \leftrightarrow H(\omega)  = \frac{1}{1-ae^{-j\omega}}
    \]
  \end{itemize}
\end{frame}

\end{document}
